(1) 线性化电子运动方程为
\[
 -i\omega m_e\mathbf v=-e(\mathbf E+\mathbf v\times B_0\hat{\mathbf z}).
\]
由于 \((\mathbf v\times B_0\hat{\mathbf z})_\pm=\mp iB_0v_\pm\)，
两个圆偏振分量解耦：
\[
 \boxed{v_\pm=-\frac{ie}{m_e(\omega\pm\omega_c)}E_\pm},\qquad
 \boxed{J_\pm=-n_e e v_\pm=
 \frac{i n_e e^2}{m_e(\omega\pm\omega_c)}E_\pm}.
\]
在本题 \(e^{-i\omega t}\) 约定下，横波 Maxwell 方程给出
\(k^2\mathbf E=\omega^2\mu_0\epsilon_0\mathbf E
+i\omega\mu_0\mathbf J\)，所以
\[
 \boxed{n_\pm^2=\frac{c^2k_\pm^2}{\omega^2}
 =1-\frac{\omega_p^2}{\omega(\omega\pm\omega_c)}}.
\]
两种偏振的正负号依赖于所定义的 \(E_x\pm iE_y\) 与时间相位，
不能直接套用另一套右旋、左旋命名而不核对符号。
当 \(B_0\to0\) 时两支均化为 \(n^2=1-\omega_p^2/\omega^2\)。

(2) 对 \(+\) 支，\(n_+^2=0\) 等价于
\(\omega(\omega+\omega_c)=\omega_p^2\)，
唯一正根为
\[
 \boxed{\omega_{\mathrm{cut},+}
 =\frac{\sqrt{\omega_c^2+4\omega_p^2}-\omega_c}{2}}.
\]
当 \(0<\omega<\omega_{\mathrm{cut},+}\) 时该支为倏逝解，
更高频率才可传播。
对 \(-\) 支，\(n_-^2=0\) 等价于
\(\omega(\omega-\omega_c)=\omega_p^2\)，正根为
\[
 \boxed{\omega_{\mathrm{cut},-}
 =\frac{\sqrt{\omega_c^2+4\omega_p^2}+\omega_c}{2}}.
\]
\(0<\omega<\omega_c\) 时分母 \(\omega(\omega-\omega_c)<0\)，
故 \(n_-^2>1\)，数学上可传播；当
\(\omega\to\omega_c^-\) 时 \(n_-^2\to+\infty\)，
而 \(\omega\to\omega_c^+\) 时 \(n_-^2\to-\infty\)，
\(\omega_c<\omega<\omega_{\mathrm{cut},-}\) 是停止带。
恰在 \(\omega_c\) 无碰撞冷流体的线性响应出现极点，
真实介质的碰撞、热运动和回旋吸收会改变奇异行为，
因此不可将无限折射率当作有限能量的物理预测。

(3) 线偏振入射满足 \(E_+(0)=E_-(0)=E_0\)。
传播 \(L\) 后
\(E_\pm(L)=E_0e^{ik_\pm L}\)，于是
\[
 E_x(L)=E_0e^{i(k_++k_-)L/2}
 \cos\frac{(k_+-k_-)L}{2},\quad
 E_y(L)=E_0e^{i(k_++k_-)L/2}
 \sin\frac{(k_+-k_-)L}{2}.
\]
因此旋转角（模 \(\pi\)）为
\[
 \boxed{\Theta=\frac{L}{2}(k_+-k_-)
 =\frac{\omega L}{2c}(n_+-n_-)}.
\]
若 \(\omega\gg\omega_p,\omega_c\)，展开平方根得到
\[
 n_\pm=1-\frac{\omega_p^2}{2\omega^2}
 \pm\frac{\omega_p^2\omega_c}{2\omega^3}
 +O\!\left(\frac{\omega_p^4}{\omega^4},
 \frac{\omega_p^2\omega_c^2}{\omega^4}\right).
\]
相减后共同项抵消，
\[
 \boxed{\Theta\simeq\frac{\omega_p^2\omega_cL}{2c\omega^2}
 =\frac{n_e e^3 B_0L}{2\epsilon_0m_e^2c\,\omega^2}}.
\]
在本题传播方向、磁场方向及 \(E_\pm\) 约定下 \(\Theta>0\)；
反向磁场使旋转符号反转。该效应随频率的平方反比衰减，
不同于无磁场等离子体的共同相位延迟。

(4) 令 \(n_\pm=\sqrt{1-\omega_p^2/[\omega(\omega\pm\omega_c)]}>0\)，
则
\[
 \frac{d(n_\pm^2)}{d\omega}
 =\frac{\omega_p^2(2\omega\pm\omega_c)}
 {\omega^2(\omega\pm\omega_c)^2}.
\]
由 \(k_\pm=\omega n_\pm/c\)，直接得到
\[
 \boxed{t_{g,\pm}=\frac Lc
 \left[n_\pm+
 \frac{\omega_p^2(2\omega\pm\omega_c)}
 {2n_\pm\omega(\omega\pm\omega_c)^2}\right]}.
\]
独立对(3)式求导即有精确关系
\[
 \boxed{t_{g,+}-t_{g,-}=2\frac{d\Theta}{d\omega}}.
\]
高频近似下，
\(k_\pm=\omega/c-\omega_p^2/(2c\omega)
\pm\omega_p^2\omega_c/(2c\omega^2)+\cdots\)，
因此
\[
 \boxed{t_{g,+}-t_{g,-}\simeq
 -\frac{2\omega_p^2\omega_cL}{c\omega^3}}.
\]
该负号意味着按本题符号约定 \(+\) 支在高频下群时延更小；
它不意味着能量或信息超光速。若 \(B_0=0\)，两支折射率
及群时延相等、旋转消失；若 \(n_e=0\)，
\(n_\pm=1\)、\(t_{g,\pm}=L/c\)、\(\Theta=0\)。
\(t_g\) 为时间量纲，\(\Theta\) 无量纲。
